\begin{textsample}{NEG Dim 2 – 2016\_c5 – Score -3.00 – 2016\_c5/enade\_2016\_ef\_4c62fccdfb.txt}  \label{ex:f2_neg_033}
CAPÍTULO 5 – DISTRIBUIÇÃO DOS CONCEITOS 5.1 - PANORAMA NACIONAL DA DISTRIBUIÇÃO DOS CONCEITOS O Conceito Enade é um indicador de qualidade que avalia o desempenho dos estudantes a partir dos resultados do Enade.
Os valores possíveis vão de 1 ( pior situação ) a 5 ( melhor situação ).
As Unidades de Observação com apenas um ou sem nenhum concluinte participante não obtêm o Conceito Enade, ficando Sem Conceito ( SC ).
A Tabela 5.1 apresenta a quantidade e distribuição de cursos de Educação Física ( Bacharelado ) participantes do Enade/2016, por faixa de conceito e Grande Região. \textbf{Enfatiza} -se, mais uma vez, que a diferença entre os cursos tabulados neste capítulo e no capítulo 2 corresponde aos cursos sem conceito, ou seja àqueles, em princípio, sem alunos concluintes que houvessem participado da prova ou àqueles que tivessem tão somente um aluno concluinte19.
Observando -se os dados da Tabela 5.1, nota -se que, dos 386 cursos participantes, 173 ( 44,8 \% ) classificaram -se com conceito 3, o valor modal.
Esse foi também o conceito modal em todas as regiões : Norte ( 50,0 \% ), Nordeste ( 50,0 \% ), Sudeste ( 42,8 \% ), Sul ( 47,3 \% ) e Centro-Oeste ( 41,9 \% ).
O conceito 4 foi o segundo mais frequente em nível nacional ( 23,3 \%, correspondendo a 90 cursos ) e o conceito 2, o terceiro ( 22,0 \%, correspondendo a 85 cursos ).
Houve, ainda, 25 cursos ( 6,5 \% ) que receberam conceito 5 e 12 cursos ( 3,1 \% ) que receberam conceito 1.
Dos 386 cursos de Educação Física ( Bacharelado ), um único curso ( 0,3 \% ) ficou sem conceito ( SC ). 19 Estes cursos com somente um concluinte não foram considerados no capítulo 2.
A região Norte participou com dez cursos ou 2,6 \% do total nacional.
Desses, cinco cursos receberam o conceito 3, o que equivale a 50,0 \% do total regional, além de ser o conceito modal na região, como já comentado.
A três cursos atribuiu -se conceitos 2, 4 e 5 ( um curso, cada ) e a dois cursos ( 20,0 \% ), conceito 1.
Nenhum curso ficou sem conceito ( SC ) nessa região.
A região Nordeste participou com 46 cursos ou 11,9 \% do total nacional.
Nessa região, 23 cursos ( 50,0 \% em termos regionais ) obtiveram conceito 3, o conceito modal para a região, como já comentado.
Os conceitos 1 e 2 foram atribuídos, respectivamente, a um e a 11 cursos ( 2,2 \% e 23,9 \%, respectivamente ).
Ainda nessa região, oito cursos ( 17,4 \% ) foram avaliados com o conceito 4, e três cursos ( 6,5 \% ), com conceito 5.
Nenhum curso ficou sem conceito ( SC ).
Dos 208 cursos participantes da região Sudeste ( 53,9 \% do total nacional ), o conceito 3 foi o conceito modal, como já anteriormente assinalado, tendo sido obtido por 89 cursos ( 42,8 \% ).
O conceito 1 foi atribuído a cinco cursos ( 2,4 \% ) e o conceito 2, a 47 cursos ( 22,6 \% ).
Receberam os conceitos 4 e 5, respectivamente, 50 cursos ( 24,0 \% ) e 17 cursos ( 8,2 \% ).
Nessa região, nenhum curso ficou sem conceito ( SC ).
Os 91 cursos da região Sul corresponderam a 23,6 \% do total nacional.
Como já indicado, houve predominância do conceito 3, atribuído a 47,3 \% dos cursos da região ( conceito modal ), ou seja, atribuído a 43 dos 91 cursos participantes na região Sul.
O conceito 1 foi atribuído a três cursos ( 3,3 \% ), o conceito 2, a 16 cursos ( 17,6 \% ), o conceito 4, a 26 cursos ( 28,6 \% ) e o conceito 5, a dois cursos ( 2,2 \% ).
Um curso ( 1,1 \% ) ficou sem conceito ( SC ) na região Sul.
Já dos 31 cursos participantes na região Centro-Oeste ( 8,0 \% do total nacional ), 13 receberam conceito 3, o conceito modal ( 41,9 \% ), como já fora destacado.
A um curso ( 3,2 \% ) foi atribuído conceito 1, e a dez ( 32,3 \% ), conceito 2.
Os conceitos restantes, 4 e 5, foram atribuídos, respectivamente, a cinco ( 16,1 \% ) e dois ( 6,5 \% ) cursos.
Nessa região, nenhum curso ficou sem conceito ( SC ).
O Gráfico 5.1 apresenta a distribuição do Conceito Enade segundo Grande Região : poligonais mais à esquerda representam Grandes Regiões com pior distribuição, e poligonais mais à direita, Grandes Regiões com melhores conceitos.
As regiões Norte ( linha azul ), Nordeste ( linha verde ) e Centro-Oeste ( linha preta ) apresentam os cursos com conceitos mais baixos, ao passo que as regiões Sudeste ( linha vermelha ) e Sul ( linha roxa ) apresentam as distribuições com valores maiores. 5.2 - CONCEITOS POR CATEGORIA ADMINISTRATIVA, POR MODALIDADE DE ENSINO E GRANDE REGIÃO A Tabela 5.2 apresenta a distribuição dos cursos participantes do Enade/2016 de Educação Física ( Bacharelado ), por Categoria Administrativa, de acordo com os conceitos por eles alcançados, segundo as Grandes Regiões brasileiras.
Dos 386 cursos participantes, 69 ( 17,9 \% ) eram ministrados em Instituições Públicas e 317 ( 82,1 \% ), em Instituições Privadas.
De acordo com as informações da Tabela 5.2, em termos nacionais, entre Instituições Públicas e Privadas, observa -se que, dos 25 cursos avaliados com conceito 5, 13 eram oferecidos em IES Públicas e 12, em IES Privadas.
Dos 69 cursos participantes de IES Públicas, o conceito 4 foi o valor modal, atribuído a 26 cursos ( 37,7 \% ).
Entre os demais cursos de Instituições Públicas participantes, quatro cursos ( 5,8 \% ) foram avaliados com conceito 1, 12 cursos ( 17,4 \% da categoria ) obtiveram conceito 2, 14 cursos ( 20,3 \% ) receberam conceito 3 e, como já comentado, 13 cursos ( 18,8 \% ) receberam conceito 5.
Nesta categoria, nenhum curso ficou sem conceito ( SC ).
Na rede privada, o conceito modal foi 3, atribuído a 159 cursos, o correspondente a 50,2 \% dos 317 cursos da categoria.
Entre os demais cursos participantes, oito ( 2,5 \% ) receberam conceito 1, e 73 ( 23,0 \% ), conceito 2.
O conceito 4 foi atribuído a 64 cursos ( 20,2 \% ) e o conceito 5, a 12 cursos ( 3,8 \% ).
Nesta Categoria Administrativa, um curso ( 0,3 \% ) ficou sem conceito ( SC ).
Os dados referentes à distribuição dos cursos segundo a Modalidade de Ensino revelam que apenas um dos 386 cursos de Educação Física ( Bacharelado ) participantes do Enade 2016 foi oferecido na modalidade Educação a Distância, equivalente 0,3 \% dos cursos.
A esse curso, foi atribuído conceito 3.
A distribuição dos conceitos nos cursos presenciais foi quase idêntica a do Brasil como um todo, considerando -se que 99,7 \% dos cursos eram desta modalidade.
Na análise por região, observa -se que, na região Norte, as Instituições Privadas participaram com seis cursos ( 60,0 \% do total regional ), a menor das proporções desta categoria entre as regiões.
Dos quais nenhum ficou sem conceito ( SC ) ou recebeu conceito 2 ou 5.
O conceito modal para as Instituições Privadas na região foi 3, atribuído a quatro cursos, correspondendo a 66,7 \% dessa categoria na região.
Outros dois cursos foram avaliados com conceitos 1 e 4 ( um curso, cada ).
As Instituições Públicas participaram com quatro cursos na região Norte ( 40,0 \% do total regional, a maior das proporções desta categoria nas regiões ), que ficaram \textbf{distribuídos} com conceitos 1, 2, 3 e 5 ( um curso, cada ).
Nenhum curso ficou sem conceito ( SC ) ou recebeu conceito 4 nessa combinação de categoria e região.
Na região Nordeste, a rede privada concentrou 33 dos 46 cursos participantes, o equivalente a 71,7 \% do total da região.
Dentre os cursos oferecidos por IES Privadas no Nordeste, 20 cursos foram avaliados com conceito 3, o conceito modal.
Um curso obteve conceito 1, nove cursos receberam conceito 2, dois ficaram com conceito 4 e um curso alcançou o conceito máximo, 5.
Nessa combinação de categoria e região, nenhum curso ficou sem conceito ( SC ).
As Instituições Públicas da região Nordeste participaram com 13 cursos ( 28,3 \% do total da região ).
Desses, seis obtiveram conceito 4, o valor modal, seguidos por outros três cursos, que receberam conceito 3.
Os conceitos 2 e 5 foram atribuídos a dois cursos, cada.
Nenhum curso ficou sem conceito ( SC ), e nenhum curso de IES Pública recebeu conceito 1 nessa região.
Na região Sudeste, a proporção de cursos da rede privada, 86,1 \%, foi mais \textbf{elevada} do que a registrada nas demais regiões brasileiras, correspondendo a 179 dos 208 cursos participantes.
Nessa categoria e região, o conceito modal foi 3, atribuído a 83 cursos.
Nenhum curso ficou sem conceito ( SC ).
Os demais foram avaliados com conceito 1 ( três cursos ), conceito 2 ( 42 cursos ), conceito 4 ( 41 cursos ) e conceito 5 ( dez cursos ).
Entre os 29 cursos oferecidos por Instituições Públicas na região Sudeste, o conceito modal foi 4, atribuído a nove cursos.
Os demais cursos receberam os conceitos 1 ( dois cursos ), 2 ( cinco cursos ), 3 ( seis cursos ) e 5 ( sete cursos ).
Na região Sudeste, nesta categoria, nenhum curso ficou sem conceito ( SC ).
As Instituições Privadas concentraram 73 dos 91 cursos participantes da região Sul, 80,2 \% do total regional.
Desses, 40 obtiveram conceito 3, o conceito modal.
Nesta combinação de Categoria Administrativa e Grande Região, os conceitos 2 e 4 foram atribuídos a 14 e a 16 cursos, respectivamente.
Dois cursos receberam conceito 1, e um ficou sem conceito ( SC ).
Além disso, nenhum curso obteve o conceito máximo 5 As Instituições Públicas na região Sul participaram com 18 cursos, 19,8 \% dos cursos da região.
O conceito modal foi 4, atribuído a dez cursos.
Aos demais cursos foram atribuídos os conceitos 1 ( um curso ), 2 ( dois cursos ), 3 ( três cursos ) e 5 ( dois cursos ). atribuídos os conceitos 1 ( um curso ), 2 ( dois cursos ), 3 ( três cursos ) e 5 ( dois cursos ).
Nenhum curso da região Sul oferecido por IES Pública ficou sem conceito ( SC ).
Na região Centro-Oeste, 26 dos 31 cursos participantes eram de Instituições Privadas ( 83,9 \% em termos regionais ).
Desses, 12 receberam conceito 3, o conceito modal.
Entre os demais, um recebeu conceito 1, oito receberam conceito 2, quatro receberam conceito 4 e um recebeu conceito 5.
Nenhum curso ficou sem conceito ( SC ).
Quanto aos cinco cursos oferecidos por Instituições Públicas na região Centro-Oeste ( 16,1 \% do total regional ), nota -se que tais cursos apresentaram o conceito 2 como modal, atribuído a dois cursos.
Três cursos receberam conceitos 3, 4 e 5 ( um curso, cada ).
Nenhum curso ficou sem conceito ( SC ) ou com conceito 1.
Como já comentado, apenas um curso foi oferecido na modalidade a distância.
Este curso estava situado na região Sudeste, correspondendo a 0,5 \% dos cursos na região.
A esse curso, mantido por IES situada na região, foi atribuído conceito 3.
O Gráfico 5.2 apresenta a distribuição do Conceito Enade segundo a Categoria Administrativa da IES.
Os cursos em IES Públicas ( linha azul ) apresentam uma poligonal mais à direita do que os em IES Privadas ( linha verde ) e, consequentemente, uma distribuição de Conceitos Enade com valores maiores. 5.3 - CONCEITOS POR ORGANIZAÇÃO ACADÊMICA E GRANDE REGIÃO Na Tabela 5.3, encontra -se a distribuição dos conceitos atribuídos aos cursos participantes do Enade/2016 na Área de Educação Física ( Bacharelado ), por Organização Acadêmica, segundo as Grandes Regiões brasileiras.
Dos 386 cursos de Educação Física ( Bacharelado ) participantes, 179 eram oferecidos em Universidades, 87 em Centros Universitários e 120 em Faculdades.
Esta distribuição corresponde a, respectivamente, 46,4 \%, 22,5 \% e 31,1 \% dos cursos.
De acordo com os dados apresentados, de todos os 25 cursos avaliados com conceito 5, 18 eram vinculados a Universidades.
Os cursos vinculados a esse tipo de Organização Acadêmica tiveram o conceito 3 como conceito modal, atribuído a 66 cursos ( 36,9 \% ).
Os demais cursos vinculados a Universidades receberam os conceitos 1 ( um curso ), 2 ( 35 cursos ), 4 ( 58 cursos ) e conceito 5 ( 18 cursos, como já mencionado ).
Apenas, um curso ficou sem conceito ( SC ).
Entre os cursos em Centros Universitários, o conceito modal foi também 3, atribuído a 43 cursos ( 49,4 \% ).
Nenhum curso vinculado a esse tipo de Organização Acadêmica ficou sem conceito ( SC ), e os outros cursos receberam os conceitos 1 ( três cursos ), 2 ( 18 cursos ), 4 ( 19 cursos ) e 5 ( quatro cursos ).
Dos 120 cursos mantidos por Faculdades, 64 ( 53,3 \% ) receberam conceito 3, o conceito modal.
O segundo conceito mais frequentemente atribuído aos cursos oferecidos por essa categoria de IES foi o conceito 2, recebido por 32 cursos.
Na sequência, 13 cursos obtiveram conceito 4 e oito cursos, conceito 1.
O conceito 5 foi atribuído a três cursos, e nenhum curso ficou sem conceito ( SC ).
Considerando -se separadamente as regiões brasileiras, verifica -se que, na região Norte, as Universidades participaram com quatro dos dez cursos participantes.
Todos eles \textbf{distribuídos} entre os conceitos 1, 2, 3 e 5 ( um curso, cada ).
Nenhum curso ficou sem conceito ( SC ) ou recebeu conceito 4 nesta combinação de categoria e região.
Os Centros Universitários da região Norte foram representados por três cursos, aos quais foram atribuídos os conceitos 1, 3 e 4 ( um curso cada ).
Nenhum curso ficou sem conceito ( SC ) ou recebeu os demais conceitos nesta combinação de categoria e região.
As Faculdades participaram com três cursos na região Norte que receberam o conceito 3.
Na região Nordeste, as Universidades participaram com 18 dos 46 cursos da Área de Educação Física ( Bacharelado ) oferecidos na região, nenhum dos quais ficou sem conceito ( SC ) ou conceito 1.
O conceito modal foi 4, atribuído a sete cursos, enquanto o conceito 5 foi atribuído a dois cursos.
Obtiveram os conceitos 2 e 3, respectivamente, três e seis cursos.
Os Centros Universitários contaram com sete cursos participantes na região Nordeste.
Esses obtiveram conceitos 2 e 4 ( um curso, cada ) e 3 ( cinco cursos, o conceito modal ).
Nenhum curso oferecido por Centros Universitários no Nordeste ficou sem conceito ( SC ) ou recebeu conceitos 1 ou 5.
As Faculdades foram representadas por 21 cursos na região Nordeste, 12 dos quais obtiveram o conceito modal 3.
O conceito 1 foi atribuído a um curso, o conceito 2, a sete cursos e o conceito 5, a um curso.
Na região Nordeste, nenhum curso oferecido por Faculdades recebeu conceito 4 ou ficou sem conceito ( SC ).
Na região Sudeste, as Universidades concentraram 91 dos 208 cursos de Educação Física ( Bacharelado ) da região.
Entre os cursos em Universidades, nessa região, o conceito modal foi 3, atribuído a 31 cursos.
Os demais cursos receberam os conceitos 2 ( 18 cursos ), 4 ( 30 cursos ) e 5 ( 12 cursos ).
Nenhum curso ficou sem conceito ( SC ) ou com conceito 1.
Os Centros Universitários participaram com 57 cursos na região Sudeste, dos quais 28 obtiveram o conceito modal 3.
O conceito 1 foi atribuído a um único curso, o conceito 2, a 15 cursos, o conceito 4, a dez cursos e, por fim, o conceito 5, a três cursos.
Nenhum curso ficou sem conceito ( SC ) nessa combinação de organização acadêmica e região.
Já as Faculdades foram representadas por 60 cursos na região Sudeste, dos quais nenhum ficou sem conceito ( SC ).
Aos demais cursos foram atribuídos os conceitos 1 ( quatro cursos ), 2 ( 14 cursos ), 3 ( 30 cursos, o conceito modal ), 4 ( dez cursos ) e 5 ( dois cursos ).
Dos 91 cursos da região Sul, 54 eram de Universidades, para os quais o conceito modal foi 3, atribuído a 24 cursos.
Nesse tipo de organização, um curso ficou sem conceito ( SC ) e nenhum curso recebeu conceito 1.
Os demais cursos receberam os conceitos 2 ( oito cursos ), 4 ( 19 cursos ) e 5 ( dois cursos ).
Dos 11 cursos participantes de Centros Universitários da região Sul, nenhum ficou sem conceito ( SC ) ou recebeu conceito 5.
O conceito modal, 4, foi atribuído a cinco cursos.
Os conceitos 1 e 2 foram atribuídos a um curso cada e o conceito 3, a quatro cursos.
Quanto aos 26 cursos vinculados a Faculdades na região Sul, 15 receberam o conceito modal 3.
O conceito 2 foi atribuído a sete cursos, e o conceitos 1 e 4, a dois cursos cada.
Nenhum curso recebeu conceito 5 ou ficou sem conceito ( SC ).
Na região Centro-Oeste, 12 dos 31 cursos eram de Universidades.
Nesse tipo de organização, nenhum curso ficou sem conceito ( SC ) ou recebeu conceito 1.
O conceito modal 2 foi atribuído a cinco cursos e o conceito 3, a quatro cursos.
Os demais conceitos, 4 e 5, foram atribuídos, respectivamente, a dois e a um curso.
Os Centros Universitários da região Centro-Oeste contaram com nove cursos, que obtiveram os conceitos 3 ( cinco cursos, conceito modal ), 4 ( dois cursos ) e dois cursos receberam os conceitos 2 ou 5 ( um curso, cada ).
Nenhum curso ficou sem conceito ( SC ) ou recebeu conceito 1.
Dos dez cursos oferecidos em Faculdades na região Centro-Oeste, nota -se que tais cursos apresentaram distribuição bimodal, uma vez que as maiores frequências se concentraram principalmente em dois conceitos, os conceitos 2 e 3, atribuídos, cada qual, a quatro cursos.
Aos demais cursos foram atribuídos os conceitos 1 e 4 ( um curso cada ).
Nenhuma Faculdade da região Centro-Oeste ficou sem conceito ( SC ) ou recebeu conceito 5.
O Gráfico 5.3 apresenta a distribuição do Conceito Enade segundo a Organização Acadêmica da IES.
Os cursos em Universidades ( linha azul ) apresentam uma poligonal mais à direita do que as demais, denotando uma melhor distribuição de Conceitos Enade.
No outro extremo, temos os cursos oferecidos em Faculdades ( linha vermelha ) com a poligonal mais à esquerda e, consequentemente, uma distribuição de Conceitos Enade com menores valores.
IES de Centros Universitários ( linha verde ) apresentam uma situação intermediária.

% matched lemmas: distribuir, elevado, enfatizar
\end{textsample}
